\backhead{OBJECTIVE TYPE QUESTIONS}{khred}
\begin{exlist}
\item Thermal time counts degrees\opt{below $T_b$}{above $T_b$}{above 30 $^\circ$C}{at night only}
\item A canopy with $kL=3$ intercepts about\opt{50}{75}{95}{100} per cent of the light
\item A $Q_{10}$ of 2 means the rate\opt{halves per 10 $^\circ$C}{doubles per 10 $^\circ$C}{doubles per 1 $^\circ$C}{is constant}
\item $K_y>1$ means the yield falls\opt{slower}{as fast as}{faster than}{independently of} the water use
\item The conversion efficiency of growth is about\opt{0.2}{0.5}{0.72}{1.0}
\end{exlist}
\objkey{1.\ (b),\quad 2.\ (c),\quad 3.\ (b),\quad 4.\ (c),\quad 5.\ (c)}
\dykband{EXERCISES IN SWAFL}
{\small Open SWAFL (desktop or online) with the benchmark project. The chapter of the \href{https://github.com/halilunaziru73-creator/Coupled-Water-Fluxes-Pore-Scale-Dynamics-Precision-Irrigation-Mulching-Feedback-Loops/raw/main/Software/v1.1.0/SWAFL_Desktop/docs/manual_v1.1.0/SWAFL_Manual_v1.1.0.pdf}{SWAFL manual} that explains each task is given in brackets; record your values in a table and compare them with what is expected (\emph{Check}).}\par\vspace{3pt}
\begin{exlist}
\item {\color{khblue}[Manual, Chapter 25]} Click every box of the crop part of the Daisy process map and list the stores and the flows of carbon.\par\noindent{\small\itshape Check: leaves, stems, storage organ and roots.}
\item {\color{khblue}[Manual, Chapter 30]} Open the dry matter of the organs of a run and read the harvest index of three years.\par\noindent{\small\itshape Check: the harvest index varies with the season.}
\item {\color{khblue}[Manual, Chapter 21]} Change the crop of the project and compare the crop water requirement.\par\noindent{\small\itshape Check: each crop has its own $K_c$ curve.}
\item {\color{khblue}[Manual, Chapter 28]} Design a Daisy experiment that varies the sowing date over four levels and run it.\par\noindent{\small\itshape Check: the yield changes with the sowing date.}
\end{exlist}
